\begin{table}[!htb]
\centering
\caption{Evaluasi Alternatif Manajemen Metadata dan \textit{Governance}}
\label{tab:metadata_comparison}
\footnotesize
\setlength{\tabcolsep}{4pt}
\begin{tabular}{|>{\raggedright\arraybackslash}m{3cm}|>{\centering\arraybackslash}m{1.5cm}|>{\centering\arraybackslash}m{2.5cm}|>{\centering\arraybackslash}m{2cm}|>{\centering\arraybackslash}m{1.5cm}|>{\centering\arraybackslash}m{1cm}|}
\hline
\multicolumn{1}{|>{\centering\arraybackslash}m{3cm}|}{\textbf{Metadata \& \textit{Governance}}} & \textbf{Maturitas} & \textbf{Ekstensibilitas} & \textbf{Komunitas} & \textbf{Integrasi} & \textbf{Total} \\
\hline
DataHub (LinkedIn) & 2 & 2 & 2 & 2 & 8 \\
\hline
Apache Atlas & 2 & 1 & 1 & 1 & 5 \\
\hline
OpenMetadata & 1 & 1 & 1 & 2 & 5 \\
\hline
Amundsen (Lyft) & 1 & 1 & 0 & 1 & 3 \\
\hline
\end{tabular}
\end{table}